\section{Pengenalan CSS dan Sintaks Dasar}

CSS (Cascading Style Sheets) adalah bahasa yang digunakan untuk mengatur tampilan dan format halaman web \cite{mdn-learn-css}. CSS memisahkan presentasi dari struktur (HTML), sehingga halaman dapat di-styling secara terpusat dan konsisten.

CSS dapat disisipkan dengan tiga cara: \textbf{inline} (atribut \texttt{style} pada elemen), \textbf{internal} (tag \texttt{<style>} di dalam \texttt{<head>}), dan \textbf{eksternal} (file \texttt{.css} terpisah yang di-link dengan \texttt{<link rel="stylesheet" href="file.css">}). Untuk proyek yang terorganisir, CSS eksternal lebih dipilih \cite{w3schools-css}.

Sintaks dasar CSS: \texttt{selector \{ properti: nilai; \}}. Selector menentukan elemen mana yang di-styling. Selector dasar meliputi: \textbf{elemen} (nama tag, misalnya \texttt{p}, \texttt{h1}), \textbf{class} (diawali titik, misalnya \texttt{.nama-class}), dan \textbf{id} (diawali \#, misalnya \texttt{\#nama-id}). Satu selector dapat memiliki banyak properti \cite{mdn-learn-css}.

\begin{lstlisting}[caption={Contoh CSS dasar}, basicstyle=\ttfamily\small, frame=single]
/* Selector elemen */
p { color: #333; font-size: 16px; }

/* Selector class */
.menu { background: #eee; padding: 10px; }

/* Selector id */
#header { border-bottom: 1px solid #ccc; }
\end{lstlisting}
